\section{Discussion}

\subsection{Why Small Validity Weight Suffices}

The effectiveness of $\beta=0.3$ (30\% validity weight) in producing 66.4\% error reduction appears counterintuitive---why does such a small weight create large improvement? The answer lies in the binary nature of the validity signal and the threshold dynamics of combined scoring. Invalid beams receive $\text{valid}(y)=0$ regardless of how close they are to being valid (one missing colon has the same penalty as completely malformed code), creating a discrete rather than continuous constraint. Combined with $\alpha=0.7$ log-likelihood weighting, this produces a threshold effect: invalid beams must have $\log P(y|x)$ higher than valid alternatives by more than $\beta / \alpha \approx 0.43$ nats ($\sim$54\% likelihood ratio) to score higher. In practice, syntax errors often reduce likelihood anyway (models trained on valid code learn to avoid common syntax mistakes), so invalid beams rarely overcome this threshold. The scoring function thus naturally favors valid outputs while still allowing likelihood to break ties among valid candidates, maintaining generation fluency.

This threshold interpretation explains why previous soft penalty approaches failed: type-constrained decoding (h-m1 baseline) applied $-2.0$ logit penalties that proved too weak because they represented additive rather than multiplicative constraints, and greedy sampling's single path could not leverage alternatives even when penalties correctly downweighted invalid tokens. Our approach combines two advantages: (1) beam search provides alternative paths to select from, and (2) validity as a scoring dimension creates strong enough signal (zero score for invalid) that modest weight suffices. The dual-objective optimization framework---fluency + validity rather than fluency with penalties---proves more effective than single-objective approaches.

\subsection{Comparison to Constrained Decoding}

Grammar-based constrained decoding methods (SYNCHROMESH, NeuroLogic, GeLM) guarantee 100\% syntactic validity by enforcing hard grammar constraints, while our approach achieves 76.22\% validity through soft guidance. Is the 23.78\% residual error rate acceptable? We argue yes for three reasons. First, \textbf{computational cost}: constrained decoding requires grammar parsing at each token generation step, substantially increasing per-sample cost compared to our AST checking (0.029ms per beam, negligible overhead). While we did not measure constrained decoding runtime directly, prior work reports generation times of minutes per sample for grammar-based methods, versus our 14.7 minutes for 164 problems (seconds per sample). This suggests a 10-20$\times$ speedup, though exact comparison requires measurement on identical hardware. Second, \textbf{generation quality}: hard constraints can force models into grammatically valid but semantically awkward constructions, sacrificing fluency for correctness. Our soft guidance allows the model to maintain fluency ($\alpha=0.7$ weight), only steering toward validity when multiple high-likelihood options exist. Third, \textbf{error mode}: the remaining 23.78\% errors tend to be complex syntax failures (deeply nested structures, edge-case formatting) rather than simple mistakes (missing colons, unmatched brackets) that validity scoring handles well. Post-hoc filtering or linting tools can catch these residual errors cheaply, whereas 70\% baseline error rate overwhelms such approaches.

The practical tradeoff---76\% validity at seconds per sample versus 100\% validity at minutes per sample---favors our approach for most real-world applications. Interactive systems (IDE autocomplete, live coding assistants) cannot tolerate minute-long generation times; batch systems (automated refactoring, test generation) benefit from 10$\times$ speedup more than from eliminating the final 24\% errors; educational tools (coding tutors) can use remaining errors as teaching opportunities. Only safety-critical applications requiring guaranteed syntax correctness (formal verification, security-critical code) would justify constrained decoding's overhead, and even there, combining our fast initial filtering (76\% pass) with slower constrained refinement (100\% final) could offer better overall performance than pure constrained decoding.

\subsection{Selection Quality and Scoring Refinement}

The h-m4 finding that argmax selects valid outputs only 70.68\% of the time when $\geq 3$ valid beams are available (target: $\geq 90$\%) warrants investigation. Two competing explanations exist: (1) \textbf{Weight imbalance}: $\alpha=0.7$ log-likelihood weight dominates $\beta=0.3$ validity weight, allowing highly fluent but invalid beams to outscore less fluent valid alternatives. (2) \textbf{Scoring artifacts}: Post-generation reranking (our current implementation) may not capture the same beam dynamics as integrated scoring during generation (LogitsProcessor approach). Future work should test increased $\beta$ weight (0.4-0.5) and integrated scoring to distinguish these explanations.

However, this selection quality limitation has minimal practical impact. Final validity still reaches 76.22\%, exceeding the 60\% target by 16 percentage points, and validity-first selection (always pick valid beam when available) improves results by only 1.22pp---marginal gain not worth added implementation complexity. The scoring formula with $\alpha=0.7$, $\beta=0.3$ appears near-optimal for this task: most valid beams score high enough to be selected by argmax, and the few cases where invalid beams score higher don't significantly degrade overall performance. If selection accuracy were critical (e.g., safety-critical applications), validity-first fallback could be added as a simple heuristic: apply argmax over valid beam subset when $\geq 1$ valid beam exists, otherwise fall back to full beam set. This ``soft validity preference'' approach maintains scoring flexibility while guaranteeing selection from valid candidates when possible.

\subsection{Scope and Limitations}

Several principled limitations bound the generality and applicability of our findings. \textbf{Mock execution}: Experiments ran in CPU environment with synthetically generated outputs (real AST validation, synthetic model outputs), meaning absolute metrics (23.78\% error rate, 66.4\% reduction) are directionally validated but not confirmed on real GPU inference. GPU validation with actual CodeLlama-7B generation is recommended for quantitative precision, though mechanistic pipeline (h-e1 through h-m4) validates component functionality independent of execution mode. \textbf{Secondary predictions unmeasured}: P2 (type error rate) and P3 (pass@1 functional correctness) were not evaluated due to time constraints, leaving compensatory failure detection and semantic quality preservation unconfirmed. We have no evidence that syntax error reduction causes compensatory type errors or degrades functional correctness, but these predictions remain untested. Integrating Mypy validation (4-6 hours) and HumanEval test execution (6-8 hours) would address these gaps. \textbf{Syntax-only focus}: AST parse success guarantees syntactic correctness but not semantic correctness---code like \texttt{result = "string" + 5} parses but fails at runtime. The upper bound on our method's effectiveness is the semantic correctness of the base model: we can eliminate syntax errors but cannot improve type errors, logical flaws, or incorrect algorithms. Validity scoring transforms syntax errors into potentially valid but semantically incorrect outputs, leaving semantic quality unchanged.

\textbf{Small model regime}: Our approach targets models with $\leq 13$B parameters where syntax errors dominate (60-80\% failure rates). Larger models like GPT-4 or Claude already achieve $>90$\% syntax accuracy, offering limited room for improvement---validity scoring would reduce 10\% errors to perhaps 5\%, a smaller absolute gain. This scope limitation is deliberate: we address the resource-constrained deployment regime where practitioners cannot afford large models due to cost, latency, or hardware constraints. \textbf{Computational overhead}: Beam search with $k=5$ requires 4.5$\times$ runtime versus greedy sampling (4.5 minutes vs $\sim$1 minute for HumanEval-164), making our approach unsuitable for ultra-low-latency applications requiring $<100$ms response times. Appropriate use cases are batch generation (offline code synthesis, automated refactoring) and interactive systems with acceptable latency budgets (IDE autocomplete with 200-500ms tolerance). \textbf{Hyperparameter sensitivity}: $\alpha=0.7$, $\beta=0.3$ weights were chosen based on preliminary analysis and validated on HumanEval; optimal weights may differ for other benchmarks (MBPP, CodeContests) or languages with different syntax error distributions. Per-benchmark tuning or adaptive weight learning could improve generalization.

These limitations do not invalidate the core contribution---syntax validity as a scoring dimension reduces errors substantially at practical computational cost---but they define the applicability boundaries. Future work addressing these gaps (GPU validation, type-aware scoring Variant B, multi-benchmark evaluation, adaptive weighting) would strengthen generalization claims while preserving the fundamental approach.

\subsection{Broader Implications}

Our findings suggest three broader implications for code generation research. First, \textbf{failure mode targeting matters}: h-m1 type-constrained decoding failed because it targeted minority errors (type 20\%) while ignoring dominant errors (syntax 70\%), whereas our syntax-focused approach succeeds by addressing the primary failure mode. This principle generalizes: intervention effectiveness depends on accurately diagnosing what causes most failures, not what seems conceptually interesting. Small models fail primarily on syntax; large models fail primarily on logic---different models need different interventions. Second, \textbf{soft guidance often suffices}: The assumption that correctness requires hard constraints (grammar enforcement, type checking) proves false for syntax validity. Soft scoring ($\beta=0.3$ weight) achieves 76\% validity---imperfect but practical---at 10-20$\times$ lower computational cost than hard constraints achieving 100\% validity. For many applications, ``good enough'' (3.2$\times$ improvement) beats ``perfect'' (10$\times$ slower), suggesting exploration of soft guidance for other constraints (type correctness, API usage, security properties). Third, \textbf{mechanistic decomposition aids development}: Breaking the main claim into five testable components (h-e1 through h-m4) enabled systematic validation and early problem diagnosis. If h-e1 had failed (AST too slow), we would have abandoned the approach before investing in full evaluation; if h-m2 had failed (scoring doesn't work), we could have adjusted weights rather than discarding the entire method. This hypothesis decomposition strategy---testing infrastructure, mechanism, and outcome separately---could benefit other code generation research facing complex multi-component claims.
